\section{Discussion and Conclusion}
\label{sec:discussion}

We have shown that sampling-based individual differential privacy is not just
\emph{vulnerable} to collusion attacks but \emph{equilibrium-incompatible}.
The pure Nash equilibria of the resulting game are intrinsically
\emph{asymmetric}---in the canonical $2$-player setting, both the naive
$(1, 8) / (8, 1)$ NE and the aware $(1, 2) / (2, 1)$ NE concentrate
vulnerability on one player while exempting the other. The budget-manipulation
attack of Kaiser et al.~\cite{kaiser2026your} is precisely the asymmetric naive
NE: a configuration that, by Proposition~\ref{prop:coalition-attack}, is the
rational best response of any coordinated subset of data holders, and that, by
Observation~\ref{obs:naive-ne}, is also reached by self-interested naive
players acting in isolation. The aware regime reduces the realised
population-level vulnerability by roughly $50\%$ but does \emph{not} restore
symmetry: it merely shifts the asymmetric equilibrium toward lower budgets.
The implication for designers of iDP systems is that the formal $(\eps_i,
\delta)$-iDP contract, while honoured at every individual record, will be
silently violated in the realised privacy of at least one of the contributors
in any pure Nash equilibrium of the strategic game it induces.

\paragraph{Limitations of this work.} Three caveats are worth noting.
First, our model assumes one-shot play. Real deployments of iDP are dynamic:
contributors join and leave; budgets can be re-negotiated; trainers
re-train. The repeated-game extension is the natural next step.
Second, we assume full information about other players' actions.
Real contracts may keep budgets private, so a Bayesian variant of the PEG
with private types is a natural and important direction. Third, our
analysis is restricted to the sampling-based realisation of iDP; the
\emph{sensitivity-based} realisation~\cite{boenisch2023have} adjusts per-sample
clipping bounds rather than per-sample sampling rates and does not exhibit the
same externality (the noise-to-sensitivity ratio is preserved by construction,
so privacy profiles are unaffected by the budget distribution). Our framework
specialises to that case as the trivial PEG in which the externality vanishes.

\paragraph{Mitigation.} The natural mitigation is to make the contract
internalise the externality. The $(\eps_i, \delta, \Delta)$-iDP formulation of
Kaiser et al.~\cite{kaiser2026your, kaissis2024beyond} adds a $\Delta$-divergence
bound on the realised vulnerability, restricting the action space of the PEG
to profiles that satisfy a per-record bound on excess advantage. We expect
this restriction to admit a unique equilibrium and to drive the Price of
Excess Vulnerability to zero (by definition). Characterising the constrained
PEG and identifying its equilibrium is a clean extension of this work.

\paragraph{Wider lesson.} The wider lesson of this paper is that an
\emph{individual} privacy contract is incomplete: it does not pin down the
collective dynamics. Where individuals interact strategically (in a single
machine-learning training, in a federated round, in any future composition of
the iDP mechanism), the contract must also constrain their joint behaviour or
the equilibrium will quietly violate every participant's expectations while
preserving every formal guarantee. The Privacy Externality Game makes this
gap explicit, formalises the equilibrium, and clarifies what an iDP contract
must do to close it.

\paragraph{Conclusion.} The promise of iDP---that every data contributor
controls her own privacy---is broken by the privacy externality of sampling-
based mechanisms. We formalised the strategic situation faced by iDP data
contributors as the Privacy Externality Game, proved monotonicity of the
externality and characterised the equilibria. The result is that the
collusion attack on iDP is not adversarial but \emph{rational}, and the
remedy must operate at the level of the contract.
